\chapter{Gauge Invariance, Cyclotron Motion, and Poisson Brackets}

A charged particle in an electromagnetic field is a small system with an
unusually rich canonical structure. Newton's equation is simple, but the
momentum that enters Hamiltonian mechanics is not merely \(m\bm v\). The
action requires a vector potential that is not unique. A uniform magnetic
field turns translation invariance into a statement about the center of a
cyclotron orbit, and an added electric field produces a steady drift
perpendicular to the applied force.

We shall develop those facts in the order in which they arise. The final
part of the chapter changes language once more, from Hamilton's equations to
Poisson brackets. That algebraic formulation is more abstract than Newton's
law, but it is the form that most clearly anticipates quantum mechanics.

\section{Four Forms of the Same Mechanics}

For the present discussion the electromagnetic field is prescribed. It may
vary from place to place, but it does not depend explicitly on time, and we
do not yet solve equations for the field itself. The only dynamical object is
a particle of mass \(m\) and charge \(q\).

It is useful to place four familiar descriptions of mechanics side by side:
\begin{enumerate}
\item Newton's equations of motion;
\item the principle of stationary action;
\item the Euler--Lagrange equations;
\item Hamilton's equations in phase space.
\end{enumerate}
They are equivalent when each is available, but they expose different
features of a theory. An equation such as \(F=ma\) tells us how a trajectory
moves. It does not, by its form alone, tell us why energy or momentum is
conserved. The action is deeper, but one normally uses it by taking a
variation and obtaining differential equations. The Lagrangian is powerful
and direct. The Hamiltonian is often the least intuitive, yet it displays the
flow of whole ensembles through phase space, makes information conservation
visible, and leads naturally toward quantum mechanics.

The charged particle will make both the power and the strangeness of these
descriptions concrete. We restrict the motion to the \(x\)-\(y\) plane. The
magnetic field is perpendicular to that plane, while an electric field may
lie within it.

\section{The Lorentz Force and the Potentials}

The nonrelativistic equation of motion is
\begin{equation}
 m\ddot{\bm x}=q\bm E+q\,\dot{\bm x}\times\bm B.
 \label{eq:l08-lorentz-force}
\end{equation}
We use an electromagnetic convention in which no explicit factor of \(c\)
appears. Other conventions insert \(1/c\) in the magnetic term; that change
does not alter the argument.

The magnetic force depends on velocity, but it is perpendicular to velocity:
\begin{equation}
 \bm v\cdot(\bm v\times\bm B)=0.
 \label{eq:l08-magnetic-power}
\end{equation}
Consequently
\begin{equation}
 \frac{d}{dt}\left(\frac12m v^2\right)
 =\bm v\cdot q(\bm v\times\bm B)=0.
 \label{eq:l08-magnetic-no-work}
\end{equation}
A magnetic field bends a trajectory without changing its speed. This is the
precise content of the statement that magnetic fields do no work. It also
explains why planar motion remains planar: with \(\bm v\) in the plane and
\(\bm B\) perpendicular to it, \(\bm v\times\bm B\) lies in the plane.

To pass from the force law to an action, we need potentials. For a region in
which the magnetic field has no monopole source,
\begin{equation}
 \bm\nabla\cdot\bm B=0,
\end{equation}
we may write, at least locally,
\begin{equation}
 \boxed{\bm B=\bm\nabla\times\bm A.}
 \label{eq:l08-vector-potential}
\end{equation}
The local qualification matters in topologically nontrivial settings and in
theories with magnetic monopoles. For the ordinary fields considered here,
one vector potential is sufficient.

\begin{classroomqa}
\audiencequestion{Why does writing \(\bm B\) as a curl imply that magnetic
charge is absent?\lecturetimestamp{00:11:34--00:13:18}}
\lecturerresponse{The divergence of every curl vanishes:
\(\bm\nabla\cdot(\bm\nabla\times\bm A)=0\). Electric field lines may begin or
end on electric charge because \(\bm\nabla\cdot\bm E\) is proportional to
charge density. In the classical electromagnetic theory being used here,
there is no corresponding magnetic charge density. Magnetic monopoles are a
possible extension of the theory, but none is part of this problem.}
\end{classroomqa}

The electrostatic field is represented by a scalar potential,
\begin{equation}
 \bm E=-\bm\nabla V,
 \qquad U=qV.
 \label{eq:l08-electric-potential}
\end{equation}
Thus \(V\) is energy per unit charge, while \(U\) is the potential energy of
this particular particle. Doubling \(q\) doubles both the electric force and
the potential energy in a fixed external field.

\begin{figure}[tbp]
\centering
\includegraphics[width=0.78\textwidth]{lecture_08_figure_01.png}
\caption{The vector potential, the electric potential, and the force law at
00:17:19. The action requires \(\bm A\), although the Newtonian equation can
be written entirely in terms of \(\bm B\).}
\label{fig:l08-potentials-board}
\end{figure}

\begin{classroomqa}
\audiencequestion{What does the symbol \(V\) measure? Is it already the
particle's potential energy?\lecturetimestamp{00:15:43--00:16:31}}
\lecturerresponse{\(V\) is the electric potential, measured in volts, or
energy per unit charge. The particle's potential energy is \(U=qV\). This
distinction keeps the prescribed field separate from the charge used to
probe it.}
\end{classroomqa}

\section{The Electromagnetic Action and Gauge Redundancy}

For planar motion, the action is
\begin{equation}
 S=\int \frac12m(\dot x^2+\dot y^2)\,dt
   +q\int_\gamma A_i\,dx_i
   -q\int V\,dt.
 \label{eq:l08-action-line-integral}
\end{equation}
The middle term is a line integral along the path \(\gamma\). Divide and
multiply its infinitesimal displacement by \(dt\):
\begin{equation}
 A_i\,dx_i=A_i\dot x_i\,dt=\bm A\cdot\bm v\,dt.
\end{equation}
The action therefore has the ordinary form \(S=\int L\,dt\), with
\begin{equation}
 \boxed{
 L=\frac12m(\dot x^2+\dot y^2)
   +q(A_x\dot x+A_y\dot y)-qV(x,y).}
 \label{eq:l08-electromagnetic-lagrangian}
\end{equation}

\begin{figure}[tbp]
\centering
\includegraphics[width=0.82\textwidth]{lecture_08_figure_02.png}
\caption{The electromagnetic action at 00:22:37. The magnetic coupling is a
line integral \(q\int\bm A\cdot d\bm x\), which becomes
\(q\bm A\cdot\bm v\) in the Lagrangian.}
\label{fig:l08-action-board}
\end{figure}

The combination \(A_i\,dx_i-V\,dt\) hints at a four-dimensional
electromagnetic potential. We do not need special relativity here, but the
formula is already arranging the scalar and vector potentials into the form
they will take in space-time.

\begin{classroomqa}
\audiencequestion{Is the repeated index in \(A_i\,dx_i\) a sum, and is the
same true of the scalar-potential term?\lecturetimestamp{00:23:36--00:24:49}}
\lecturerresponse{Yes for the vector term:
\(A_i\,dx_i=A_x\,dx+A_y\,dy\) in the plane. The quantity \(V\) is a scalar,
so there is no component sum in \(V\,dt\). To avoid confusing electric
potential \(V\) with velocity, we shall write the latter as \(\bm v\).}
\end{classroomqa}

The action has introduced more description than the force law seemed to
need. The reason is that \(\bm A\) is not unique. For any smooth scalar
function \(\lambda(\bm x)\),
\begin{equation}
 \bm A\longrightarrow\bm A'=\bm A+\bm\nabla\lambda
 \label{eq:l08-gauge-transformation}
\end{equation}
leaves the magnetic field unchanged, since mixed partial derivatives
commute:
\begin{equation}
 \bm\nabla\times\bm A'
 =\bm\nabla\times\bm A
  +\bm\nabla\times\bm\nabla\lambda
 =\bm B.
\end{equation}
In the time-independent setting, the Lagrangian changes by a total
derivative,
\begin{equation}
 L' = L+q\bm\nabla\lambda\cdot\dot{\bm x}
    = L+q\frac{d\lambda}{dt}.
 \label{eq:l08-gauge-total-derivative}
\end{equation}
Hence
\begin{equation}
 S'=S+q\lambda(\bm x_f)-q\lambda(\bm x_i).
\end{equation}
The endpoint contribution does not affect a variation with fixed endpoints.
The Lagrangian changes, but the trajectory does not. This is gauge
invariance: different mathematical descriptions represent the same physical
field.

The redundancy is not a defect that can simply be discarded. We could stay
with the Lorentz-force equation and never mention \(\bm A\), but the local
action, canonical momentum, and Hamiltonian all require it. Modern physics
often reaches its most useful formulations by introducing such redundant
variables and then identifying the transformations that leave observables
unchanged.

\section{Recovering the Lorentz Force}

We should not accept the Lagrangian merely because it has the right-looking
symbols. Let us derive the \(x\)-component of the equation of motion. First,
\begin{equation}
 p_x=\frac{\partial L}{\partial\dot x}
    =m\dot x+qA_x.
 \label{eq:l08-canonical-px}
\end{equation}
This is our first warning that canonical momentum and mechanical momentum
are different:
\begin{equation}
 \bm p_{\mathrm{can}}=m\bm v+q\bm A,
 \qquad
 \bm p_{\mathrm{mech}}=m\bm v.
 \label{eq:l08-two-momenta}
\end{equation}
The canonical momentum changes when \(\bm A\) changes gauge. It is a phase
space coordinate adapted to the chosen Lagrangian, not by itself a
gauge-invariant observable.

The Euler--Lagrange equation gives
\begin{equation}
 \frac{d}{dt}\frac{\partial L}{\partial\dot x}
 =\frac{\partial L}{\partial x}.
\end{equation}
Although \(A_x\) has no explicit time dependence, its value at the moving
particle changes. The chain rule gives
\begin{equation}
 \frac{dA_x}{dt}
 =\frac{\partial A_x}{\partial x}\dot x
 +\frac{\partial A_x}{\partial y}\dot y.
 \label{eq:l08-chain-rule-Ax}
\end{equation}
Therefore
\begin{align}
 m\ddot x
 +q\frac{\partial A_x}{\partial x}\dot x
 +q\frac{\partial A_x}{\partial y}\dot y
 &=q\frac{\partial A_x}{\partial x}\dot x
  +q\frac{\partial A_y}{\partial x}\dot y
  -q\frac{\partial V}{\partial x}.
\end{align}
The common terms cancel, leaving
\begin{align}
 m\ddot x
 &=q\left(
 \frac{\partial A_y}{\partial x}
 -\frac{\partial A_x}{\partial y}
 \right)\dot y-q\frac{\partial V}{\partial x}\notag\\
 &=qB_z\dot y+qE_x.
 \label{eq:l08-recovered-lorentz-x}
\end{align}

\begin{figure}[tbp]
\centering
\includegraphics[width=0.72\textwidth]{lecture_08_figure_03.png}
\caption{The surviving curl combination at 00:38:20. The separate
derivatives of \(\bm A\) combine into the gauge-invariant field \(B_z\).}
\label{fig:l08-lorentz-recovery-board}
\end{figure}

The vector potential has disappeared from the final equation except through
its curl. Thus the Lagrangian uses a redundant variable, while its
Euler--Lagrange equation returns the gauge-invariant force law.

\section{A Uniform Magnetic Field}

We now use the canonical formalism rather than merely checking it. Fix
ordinary right-handed axes with \(+\hat{\bm z}\) pointing out of the board,
and write the magnetic field into the board as
\begin{equation}
 \bm B=-B\hat{\bm z},\qquad B>0.
 \label{eq:l08-field-convention}
\end{equation}
Then \(B_z=-B\). Two convenient vector potentials for this same field are
\begin{align}
 \text{gauge I:}&\qquad A_x=By, & A_y&=0,
 \label{eq:l08-gauge-one}\\
 \text{gauge II:}&\qquad A_x=0, & A_y&=-Bx.
 \label{eq:l08-gauge-two}
\end{align}
Indeed, each gives
\(\partial_xA_y-\partial_yA_x=-B\). A symmetric choice,
\begin{equation}
 A_x=\frac{B}{2}y,
 \qquad A_y=-\frac{B}{2}x,
\end{equation}
does the same. No experiment can distinguish these gauges.

Set the electric field to zero and begin with gauge I. The Lagrangian is
\begin{equation}
 L_1=\frac12m(\dot x^2+\dot y^2)+qBy\dot x,
 \label{eq:l08-L-gauge-one}
\end{equation}
and the canonical momenta are
\begin{equation}
 p_x=m\dot x+qBy,
 \qquad p_y=m\dot y.
 \label{eq:l08-momenta-gauge-one}
\end{equation}

\begin{figure}[tbp]
\centering
\includegraphics[width=0.8\textwidth]{lecture_08_figure_05.png}
\caption{Gauge I at 00:46:20. The Lagrangian contains \(qBy\dot x\), so the
canonical momentum is \(p_x=m\dot x+qBy\).}
\label{fig:l08-gauge-one-board}
\end{figure}

\begin{classroomqa}
\audiencequestion{Which of these two canonical momenta is conserved?
\lecturetimestamp{00:46:18--00:48:38}}
\lecturerresponse{The Lagrangian depends on \(y\) but not on \(x\). A
translation \(x\mapsto x+\epsilon\) changes no term appearing in \(L_1\), so
\(x\) is cyclic and \(p_x\) is conserved.
A translation in \(y\) changes the term \(qBy\dot x\), so \(p_y=m\dot y\) is
not the conserved canonical quantity in this gauge.}
\end{classroomqa}

Choose the conserved value \(p_x=0\) for the moment. Then
\begin{equation}
 \boxed{\dot x=-\frac{qB}{m}y.}
 \label{eq:l08-xdot}
\end{equation}

\begin{figure}[tbp]
\centering
\includegraphics[width=0.7\textwidth]{lecture_08_figure_06.png}
\caption{The first velocity relation at 00:50:10, obtained from the conserved
canonical momentum in gauge I.}
\label{fig:l08-xdot-board}
\end{figure}

Now use gauge II. Its Lagrangian is
\begin{equation}
 L_2=\frac12m(\dot x^2+\dot y^2)-qBx\dot y,
\end{equation}
and \(y\) is cyclic. Thus
\begin{equation}
 p_y=m\dot y-qBx
\end{equation}
is conserved. Choosing \(p_y=0\) gives
\begin{equation}
 \boxed{\dot y=\frac{qB}{m}x.}
 \label{eq:l08-ydot}
\end{equation}

\begin{figure}[tbp]
\centering
\includegraphics[width=0.7\textwidth]{lecture_08_figure_07.png}
\caption{The two velocity relations together at 00:52:00. Their form already
identifies uniform circular motion.}
\label{fig:l08-two-velocities-board}
\end{figure}

It may seem peculiar to use one gauge to expose the \(x\)-relation and
another to expose the \(y\)-relation. What is being combined is not two
incompatible descriptions at one instant, but two gauge-independent
statements about the same physical orbit. Each gauge makes a different
translation symmetry manifest.

\subsection{Cyclotron motion and the charge-to-mass ratio}

The pair of first-order equations is solved by
\begin{equation}
 x=r\cos(\omega_c t),
 \qquad y=r\sin(\omega_c t),
\end{equation}
because
\begin{equation}
 \dot x=-\omega_c y,
 \qquad \dot y=\omega_c x.
\end{equation}
Comparison with Eqs.~\eqref{eq:l08-xdot} and \eqref{eq:l08-ydot} gives the
signed cyclotron frequency
\begin{equation}
 \boxed{\omega_c=\frac{qB}{m}},
 \qquad |\omega_c|=\frac{|q|B}{m}.
 \label{eq:l08-cyclotron-frequency}
\end{equation}
The sign fixes the direction of rotation; the magnitude fixes the number of
radians swept out per unit time. A larger charge bends the trajectory more
rapidly, while a larger mass bends it more slowly.

Historically, magnetic deflection experiments naturally measured \(q/m\).
The trajectory alone did not separately reveal the electron's charge and
mass; another kind of measurement was needed to disentangle them.

\begin{classroomqa}
\audiencequestion{Does the sign of the cyclotron frequency matter?
\lecturetimestamp{00:54:19--00:55:24}}
\lecturerresponse{Its magnitude is \(|q|B/m\). The sign records orientation.
Reversing the charge or the magnetic field reverses the direction in which
the orbit is traversed, without changing its radius or speed.}
\end{classroomqa}

\subsection{What the conserved momenta measure}

A uniform field does not privilege the origin, so the circle may be centered
at any fixed point \((x_0,y_0)\):
\begin{equation}
 x=x_0+r\cos(\omega_c t),
 \qquad y=y_0+r\sin(\omega_c t).
 \label{eq:l08-shifted-orbit}
\end{equation}
Substitute this orbit into the two conserved canonical momenta. In gauge I,
\begin{align}
 p_x
 &=m\dot x+qBy\notag\\
 &=-m\omega_c r\sin(\omega_c t)
   +qB\bigl[y_0+r\sin(\omega_c t)\bigr]\notag\\
 &=qB y_0.
 \label{eq:l08-px-center}
\end{align}
In gauge II,
\begin{align}
 p_y
 &=m\dot y-qBx\notag\\
 &=m\omega_c r\cos(\omega_c t)
   -qB\bigl[x_0+r\cos(\omega_c t)\bigr]\notag\\
 &=-qB x_0.
 \label{eq:l08-py-center}
\end{align}
The oscillating terms cancel because \(m\omega_c=qB\). The strange conserved
momenta have a simple physical meaning: they encode the fixed guiding center
of the orbit.

\begin{figure}[tbp]
\centering
\includegraphics[width=0.8\textwidth]{lecture_08_figure_08.png}
\caption{The shifted circular orbit and its conserved canonical momenta at
01:01:20. The board gives \(p_x=qBy_0\) and \(p_y=-qBx_0\).}
\label{fig:l08-guiding-center-board}
\end{figure}

\begin{figure}[tbp]
\centering
\begin{tikzpicture}[>=Stealth,scale=0.9]
  \draw[->] (-0.6,0) -- (6.2,0) node[right] {\(x\)};
  \draw[->] (0,-0.6) -- (0,4.7) node[above] {\(y\)};
  \coordinate (C) at (3.7,2.5);
  \draw[thick,blue!60!black] (C) circle (1.35);
  \fill (C) circle (1.8pt) node[above right] {\((x_0,y_0)\)};
  \draw[dashed] (3.7,0) -- (C);
  \draw[dashed] (0,2.5) -- (C);
  \draw[->,red!65!black,thick]
    ($(C)+(35:1.35)$) arc[start angle=35,end angle=105,radius=1.35];
  \node[blue!60!black] at (5.3,3.55) {cyclotron orbit};
  \node at (3.7,-0.32) {\(x_0\)};
  \node[left] at (-0.05,2.5) {\(y_0\)};
  \foreach \x in {0.8,1.6,2.4,3.2,4.0,4.8,5.6}
    \node at (\x,4.25) {\(\scriptstyle\otimes\)};
  \node[anchor=west] at (0.45,4.25) {\(\bm B\) into the page};
\end{tikzpicture}
\caption{Reconstruction of the geometry used at 00:57:52--01:01:35. The
canonical constants identify the orbit center rather than the instantaneous
mechanical momentum.}
\label{fig:l08-guiding-center-reconstruction}
\end{figure}

\begin{classroomqa}
\audiencequestion{Are the dimensions of \(qBy_0\) really those of momentum?
\lecturetimestamp{01:02:19--01:02:45}}
\lecturerresponse{Yes, provided one uses one electromagnetic unit convention
consistently. In the convention of the present Lagrangian, \(q\bm A\) has the
dimensions of momentum. Since a uniform-field potential scales as
\(A\sim B\times\text{length}\), the product \(qBy_0\) has the same dimensions
as \(m\dot x\).}
\end{classroomqa}

\section{Electric Drift and the Hall Response}

Restore a constant electric field in the \(+x\)-direction. Choose
\begin{equation}
 V(x)=-Ex,
 \qquad \bm E=E\hat{\bm x}.
\end{equation}
With the field convention in Eq.~\eqref{eq:l08-field-convention}, the
component equations are
\begin{align}
 m\ddot x+qB\dot y&=qE,
 \label{eq:l08-crossed-field-x}\\
 m\ddot y-qB\dot x&=0.
 \label{eq:l08-crossed-field-y}
\end{align}
The electric potential depends on \(x\), so \(x\)-translation invariance is
lost and \(p_x\) is no longer conserved:
\begin{equation}
 \dot p_x=\frac{\partial L}{\partial x}=qE.
\end{equation}

Rather than solve the most general orbit, look first for a trajectory with no
acceleration. Setting \(\ddot x=\ddot y=0\) in
Eqs.~\eqref{eq:l08-crossed-field-x}--\eqref{eq:l08-crossed-field-y} gives
\begin{equation}
 \dot x=0,
 \qquad
 \boxed{\dot y=\frac{E}{B}.}
 \label{eq:l08-hall-drift}
\end{equation}

\begin{figure}[tbp]
\centering
\includegraphics[width=0.78\textwidth]{lecture_08_figure_09.png}
\caption{The no-acceleration solution at 01:08:20. The electric force points
along \(x\), while the steady drift has \(\dot y=E/B\).}
\label{fig:l08-hall-board}
\end{figure}

The coordinate-free form is
\begin{equation}
 \boxed{\bm v_D=\frac{\bm E\times\bm B}{B^2}.}
 \label{eq:l08-exb-drift}
\end{equation}
For \(\bm E=E\hat{\bm x}\) and
\(\bm B=-B\hat{\bm z}\), this vector points along \(+\hat{\bm y}\). The
charge cancels: particles of different charge magnitude share the same
\(\bm E\times\bm B\) drift. The full crossed-field orbit is a cyclotron
motion whose guiding center moves with \(\bm v_D\).

With only a magnetic field, the no-acceleration solution is the zero-radius
orbit: the particle simply sits still. Turning on \(\bm E\) changes that
solution into steady transverse motion. A distribution of charges therefore
develops a sideways flow. This is the elementary mechanical structure behind
the Hall response.

\begin{classroomqa}
\audiencequestion{Why does a force applied along \(x\) produce motion along
\(y\)?\lecturetimestamp{01:07:07--01:10:45}}
\lecturerresponse{A trial velocity along \(y\) generates a magnetic force
along \(x\). At the special speed \(E/B\), that magnetic force exactly
cancels the electric force. The net force and acceleration vanish, so the
transverse velocity persists.}
\end{classroomqa}

\begin{classroomqa}
\audiencequestion{Is this related to Hall-effect devices used with motors or
disk drives?\lecturetimestamp{01:10:47--01:11:38}}
\lecturerresponse{Hall sensors do exploit transverse electromagnetic
response and are commonly used to sense magnetic fields, position, or rotor
commutation. The detailed engineering of a motor is beyond this mechanics
problem, so we should not infer a particular mechanism without examining the
device.}
\end{classroomqa}

The same kind of perpendicular response appears in gyroscopic precession. A
torque that one might naively expect to tip a spinning body instead produces
motion around a transverse direction. The physical systems differ, but their
coupled first-order equations have a closely related structure.

\subsection{Completing the Hamiltonian exercise}

The charged-particle Hamiltonian follows directly from the exercise posed in
the lecture. From
\begin{equation}
 \bm p=m\bm v+q\bm A
\end{equation}
we have
\begin{equation}
 \bm v=\frac{\bm p-q\bm A}{m}.
\end{equation}
The Legendre transform gives
\begin{align}
 H&=\bm p\cdot\bm v-L\notag\\
  &=\frac{(\bm p-q\bm A)^2}{2m}+qV.
 \label{eq:l08-charged-hamiltonian}
\end{align}
Thus kinetic energy depends on the gauge-invariant mechanical combination
\(\bm p-q\bm A=m\bm v\), even though Hamilton's equations use the canonical
coordinate \(\bm p\). Applying those equations reproduces the Lorentz force,
as it must.

\section{Poisson Brackets}

We now leave the charged-particle example and abstract the structure of
Hamiltonian mechanics. For canonical coordinates \((q_i,p_i)\), Hamilton's
equations are
\begin{equation}
 \dot q_i=\frac{\partial H}{\partial p_i},
 \qquad
 \dot p_i=-\frac{\partial H}{\partial q_i}.
 \label{eq:l08-hamilton-equations}
\end{equation}
There are twice as many first-order equations as there are configuration
coordinates, in place of the same number of second-order Lagrange equations.

Let \(A(q_i,p_i)\) be any function on phase space with no explicit time
dependence. Along a trajectory,
\begin{align}
 \dot A
 &=\sum_i\left(
 \frac{\partial A}{\partial q_i}\dot q_i
 +\frac{\partial A}{\partial p_i}\dot p_i
 \right)\notag\\
 &=\sum_i\left(
 \frac{\partial A}{\partial q_i}\frac{\partial H}{\partial p_i}
 -\frac{\partial A}{\partial p_i}\frac{\partial H}{\partial q_i}
 \right).
 \label{eq:l08-Adot-chain-rule}
\end{align}
This combination motivates a definition. For any two phase-space functions
\(A\) and \(B\), their Poisson bracket is
\begin{equation}
 \boxed{
 \{A,B\}=\sum_i\left(
 \frac{\partial A}{\partial q_i}\frac{\partial B}{\partial p_i}
 -\frac{\partial A}{\partial p_i}\frac{\partial B}{\partial q_i}
 \right).}
 \label{eq:l08-poisson-definition}
\end{equation}

\begin{figure}[tbp]
\centering
\includegraphics[width=0.8\textwidth]{lecture_08_figure_10.png}
\caption{The Poisson-bracket definition at 01:20:40. The relative minus sign
is the source of its antisymmetry.}
\label{fig:l08-poisson-definition-board}
\end{figure}

Equation~\eqref{eq:l08-Adot-chain-rule} is now simply
\begin{equation}
 \boxed{\dot A=\{A,H\}.}
 \label{eq:l08-hamilton-evolution-bracket}
\end{equation}
If \(A\) has explicit time dependence, the complete statement is
\begin{equation}
 \frac{dA}{dt}=\frac{\partial A}{\partial t}+\{A,H\}.
\end{equation}

The bracket formula contains both Hamilton equations. With one canonical
pair \((Q,P)\),
\begin{align}
 \dot Q&=\{Q,H\}=\frac{\partial H}{\partial P},\\
 \dot P&=\{P,H\}=-\frac{\partial H}{\partial Q}.
\end{align}
So Poisson brackets do not introduce new dynamics. They compress the old
dynamics into an algebraic operation on functions.

\begin{classroomqa}
\audiencequestion{What happens if \(A\) is the constant function \(1\), or
if it is simply \(P\) or \(Q\)?\lecturetimestamp{01:22:28--01:25:34}}
\lecturerresponse{A constant has zero derivatives, so
\(\{1,H\}=0\). Choosing \(A=P\) gives
\(\dot P=-\partial H/\partial Q\); choosing \(A=Q\) gives
\(\dot Q=\partial H/\partial P\). The simplest choices recover the two
Hamilton equations directly.}
\end{classroomqa}

\subsection{The fundamental algebra}

Interchanging the two entries changes the sign:
\begin{equation}
 \{A,B\}=-\{B,A\}.
 \label{eq:l08-antisymmetry}
\end{equation}
In particular, any function has zero bracket with itself. For the canonical
variables,
\begin{equation}
 \boxed{
 \{q_i,q_j\}=0,
 \qquad \{p_i,p_j\}=0,
 \qquad \{q_i,p_j\}=\delta_{ij}.}
 \label{eq:l08-fundamental-brackets}
\end{equation}
Consequently \(\{p_i,q_j\}=-\delta_{ij}\): reversing the entries reverses the
sign.

\begin{figure}[tbp]
\centering
\includegraphics[width=0.78\textwidth]{lecture_08_figure_11.png}
\caption{Antisymmetry and the fundamental canonical brackets at 01:29:40.
The Kronecker delta is one for matching indices and zero otherwise.}
\label{fig:l08-fundamental-brackets-board}
\end{figure}

\begin{classroomqa}
\audiencequestion{Can one see immediately that \(\{P,P\}=0\), without
expanding the definition?\lecturetimestamp{01:27:06--01:28:05}}
\lecturerresponse{Yes. Antisymmetry says \(\{P,P\}=-\{P,P\}\). The only
quantity equal to its own negative is zero. The same argument gives
\(\{Q,Q\}=0\).}
\end{classroomqa}

For one canonical pair, bracketing with \(P\) or \(Q\) acts as
differentiation:
\begin{equation}
 \{P,F\}=-\frac{\partial F}{\partial Q},
 \qquad
 \{Q,F\}=\frac{\partial F}{\partial P}.
 \label{eq:l08-bracket-as-derivative}
\end{equation}
More generally,
\begin{equation}
 \{p_i,F\}=-\frac{\partial F}{\partial q_i},
 \qquad
 \{q_i,F\}=\frac{\partial F}{\partial p_i}.
\end{equation}
This is the symplectic asymmetry of classical mechanics: exchanging a
coordinate with its momentum brings in a minus sign.

The bracket is linear in each entry. For constants \(\alpha,\beta\),
\begin{equation}
 \{\alpha A+\beta C,B\}
 =\alpha\{A,B\}+\beta\{C,B\}.
 \label{eq:l08-poisson-linearity}
\end{equation}

\begin{figure}[tbp]
\centering
\includegraphics[width=0.74\textwidth]{lecture_08_figure_13.png}
\caption{Linearity of the bracket at 01:36:20. Multiplication by a constant
and addition pass through the bracket.}
\label{fig:l08-linearity-board}
\end{figure}

It also obeys the ordinary product rule in either entry. Expanding the
derivatives of \(AB\) gives
\begin{align}
 \{AB,C\}
 &=\frac{\partial(AB)}{\partial q_i}
   \frac{\partial C}{\partial p_i}
  -\frac{\partial(AB)}{\partial p_i}
   \frac{\partial C}{\partial q_i}\notag\\
 &=A\left(
   \frac{\partial B}{\partial q_i}\frac{\partial C}{\partial p_i}
  -\frac{\partial B}{\partial p_i}\frac{\partial C}{\partial q_i}
  \right)
 +B\left(
   \frac{\partial A}{\partial q_i}\frac{\partial C}{\partial p_i}
  -\frac{\partial A}{\partial p_i}\frac{\partial C}{\partial q_i}
  \right)\notag\\
 &=A\{B,C\}+B\{A,C\}.
 \label{eq:l08-poisson-product-rule}
\end{align}
Repeated indices are summed in this calculation.

\begin{figure}[!htbp]
\centering
\includegraphics[width=0.68\textwidth]{lecture_08_figure_12.png}
\caption{The product-rule expansion at 01:39:00. Grouping the terms yields
\(\{AB,C\}=A\{B,C\}+B\{A,C\}\).}
\label{fig:l08-product-rule-board}
\end{figure}

\begin{classroomqa}
\audiencequestion{Besides the fundamental brackets, what rules are needed
to calculate brackets of general functions?\lecturetimestamp{01:33:48--01:40:48}}
\lecturerresponse{Antisymmetry fixes the order, linearity handles sums and
constant multiples, and the product rule handles products. Together with
\(\{q_i,p_j\}=\delta_{ij}\), these rules determine brackets of polynomials and,
by suitable limits, the ordinary smooth functions used in mechanics.}
\end{classroomqa}

We can therefore state classical mechanics in a compact way. Begin with
canonical variables obeying the fundamental Poisson brackets, specify a
Hamiltonian, and evolve every observable by
\(\dot A=\{A,H\}\). This formulation follows from Hamilton's equations, but
it makes the algebraic structure stand on its own.

\newpage
\begin{samepage}
The reason for carrying the abstraction this far is not merely elegance. In
quantum mechanics, Poisson brackets are replaced in a precise way by
commutators, and the same algebraic pattern becomes the law of motion for
operators. Heisenberg and Dirac made that connection central. What looks
like an ornate reformulation of classical mechanics here will become one of
the shortest routes into quantum dynamics.
\end{samepage}
